%--------------------------------------------------------------------
% Pakete
%--------------------------------------------------------------------

\usepackage{epsfig, psfrag}
\usepackage[multiple]{footmisc}
\usepackage{lipsum}
\usepackage{mdframed}

\usepackage[withpage]{acronym}      % Abkürzungen verwalten
\usepackage{algorithm}
%\usepackage{algorithmicx}
\usepackage{algpseudocode}
\usepackage{amsmath}
\usepackage{amssymb}                % Erweiterter Mathematiksatz
\usepackage{amsthm}

\usepackage{booktabs}               % Schöne Tabellen mit \toprule\midrule\bottomrule
\usepackage{calc}                   % Rechnen in LaTeX
\usepackage[font=small,labelfont=it,format=hang]{caption}[2008/04/01]  % 
%\usepackage{enumitem}               % Listen flexibel anpassen, z.B. keine Einrückung
\usepackage[right]{eurosym}         % Eurosymbol mit \euro und \EUR{Zahl}
\usepackage{float}                  % Gleitobjekte fixieren
\usepackage[T1]{fontenc}            % Korrektes Trennen von Wörtern mit Umlauten
\usepackage{footnote}
\usepackage[left=3cm,right=3cm,top=3cm,bottom=3cm]{geometry}
\usepackage{graphicx}               % Grafikeinbindung
\usepackage[colorlinks=false]{hyperref}             % Links
\usepackage{hyphenat}               % Umbruch
\usepackage{ifpdf}
\usepackage{ifthen}                 % if-Verzeigung
\usepackage[utf8]{inputenc}         % Direktes Eingeben deutscher Umlaute
\usepackage{listings}               % Quellcode einbinden und formatieren
\usepackage{lmodern}
\usepackage{multirow}               % Tabellenzelle über mehrere Zeilen
\usepackage{natbib}                 % Zitieren nach [Autor, Jahr] und [1]

\usepackage{nomencl}
\usepackage{paralist}               % Kompakte Listen (Aufzählungen)
\usepackage{setspace}
\usepackage{xspace}
% Bei Verwendung von lmodern fehlen ein paar Zeichen, z.B. Punkt bei itemize
% --> Font Warning: Font shape `OMS/lmr/m/n' undefined
% --> Some font shapes were not available, defaults substituted.
% Deshalb Paket 'textcomp' laden.
\usepackage{shortvrb}
\usepackage{tabularx}               % Tabelle mit fester Gesamtbreite und 
                                    % variabler Spaltenbreite
\usepackage{textcomp}               % Zusätzliche Sonderzeichen
\usepackage[table]{xcolor}          % Erweitertes Farbpaket mit vielen Farbmodellen
\usepackage{xtab}
\usepackage[ngerman,english,german]{babel}
\usepackage{pgf-pie}
\usepackage{pst-uml}
\usepackage{pgf-umlsd}
\usepackage{pgf-umlcd}

\graphicspath{{./images}}

\usepackage{url}
\usepackage{breakurl}
